\documentclass[10pt]{article}

% Specify the margins and text size.
\setlength{\textwidth}{6.4in}
\setlength{\textheight}{9.5in}
\setlength{\oddsidemargin}{0pt}
\setlength{\evensidemargin}{0pt}
\setlength{\topmargin}{0pt}
\setlength{\hoffset}{.05in}
\setlength{\voffset}{-.8in}

\setlength{\parskip}{5pt}
\setlength{\parindent}{0pt}

% Load some fonts and symbol packages
\usepackage{latexsym}
\usepackage{pifont}       % contains 'star' symbol for counterinsurgency handbook title
\usepackage{yfonts} 
\usepackage{amsmath}
\usepackage{amsfonts}

\usepackage{graphicx}     % actually, this is loaded by pstricks
\usepackage[T1]{fontenc}
\usepackage{ifthen}
\usepackage{pstricks,pst-grad,pst-text,pst-node,multido,pst-plot,calc,pst-3dplot}
%\usepackage[all]{xy}
%\usepackage{animate}

% The hyperref package inserts links into the pdf.
\definecolor{MyLinkColor}{rgb}{.1,.2,1}
\definecolor{MyCiteColor}{rgb}{.1,1,.2}
\definecolor{MyURLColor}{rgb}{.2,.2,.6}
\usepackage[backref=true,pagebackref=false,hyperindex,colorlinks=true,
  linkcolor=MyLinkColor,urlcolor=MyURLColor]{hyperref}


% The tweaklist package is something I found on the web.  It provides a simple interface
% for making changes to spacing used in the itemize and enumerate environments.  Comment
% this out if you don't care to use tweaklist.
\usepackage{tweaklist}
\renewcommand{\itemhook}{\setlength{\parskip}{2pt}\setlength{\parsep}%
{1pt}\setlength{\topsep}{0pt}\setlength{\itemsep}{0pt}}

\newcommand{\U}{\underline{\hspace{5pt}}}

\usepackage{listings}
\newcounter{EX}\setcounter{EX}{1}
\newcommand{\EXERCISE}{\arabic{EX}.\stepcounter{EX} }

\begin{document}
\pagestyle{empty}
\lstset{language=R, showspaces=false, showstringspaces=false}

\href{http://www.su.edu}{\includegraphics[height=1.75cm]{sulogo.eps}}
\vspace{-1.69cm}

{\small \hfill
\begin{tabular}{cl}
& Math 207\\& Introduction to Statistics\\
\end{tabular}
}
\setlength{\baselineskip}{1.05\baselineskip}
\medskip
\medskip

\begin{center}
\textbf{\large  First Experiments with the NHANESIII Data Set}
\end{center}

\newcommand{\SUBX}{\smallskip\hspace{10pt}}
\newcommand{\BSK}{\vspace{.14in}}

For much of this course we will explore the NHANESIII data set
in order to illustrate the concepts that we cover.  This is a large survey of
health statistics of U.S. residents.  The data set website is:\vspace{-8pt}
\begin{center}
\href{https://wwwn.cdc.gov/nchs/nhanes/nhanes3/datafiles.aspx}{https://wwwn.cdc.gov/nchs/nhanes/nhanes3/datafiles.aspx}\vspace{-8pt}
\end{center}
The data consists of four large files.  On the CDC website these are in a format 
developed for the SAS statistical software package.  I had these converted to 
csv (comma-separated-value) text files then put them in a zip file which is on our course website.

\EXERCISE \textbf{Download the data set}.
\begin{itemize}
\item Create a folder called \texttt{math207} on your desktop.
\item Download the \texttt{NHANESIII.zip} file from Canvas.
\item Move the \texttt{NHANESIII.zip} file into your \texttt{math207} directory.
\item Open a terminal and type the following:\\
  \texttt{cd  Desktop}\\
  \texttt{cd  math207}\\
  \texttt{unzip  NHANESIII.zip}
\end{itemize}

The \texttt{math207} folder should contain the following files:\vspace{-6pt}
\begin{center}
\begin{tabular}{|ll|}\hline
Data File & Contents\\\hline
adult.csv  & survey information about adults subjects\\
youth.csv  & survey information about youth subjects\\
lab.csv    & blood and urine test data\\
exam.csv   & large collection of medical exam data\\[4pt]
%
ADULT-acc.pdf & information file for the adult.csv data file\\
YOUTH-acc.pdf & information file for the youth.csv data file\\
lab-acc.pdf   & information file for the lab.csv data file\\
exam-acc.pdf  & information file for the exam.csv data file\\\hline
\end{tabular}
\end{center}\vspace{-3pt}


\EXERCISE \textbf{Experimenting with the Adult Survey Data}. 
Open RStudio and the ADULT-acc.pdf code book (statisticians' term for a 
document that describes what is in a data set).
Try the following in RStudio.  
Explain to yourself what \texttt{R} thinks it is doing.\vspace{-6pt}
\begin{verbatim}
# First, load the data <- This is a comment, you don't have to type it.
# USE YOUR USERNAME not the word "username".
adult = read.csv("/Users/username/Desktop/math207/adult.csv", header=TRUE)
class(adult)

# This will tell you the number of rows (20,050) and columns (1238).
dim(adult)

# Here are the 1238 (cryptic) variable names.  They are desribed briefly
# on pages 30 -- 66 of the code book.  Pages 66 -- 391 give some basic
# statistics for each variable.  The rest of the document gives more
# detailed descriptions.
names(adult)

# The first variable of interest is SEQN.  It stores the anonymous ID
# numbers for participants.  Each subject has data stored in multiple files.
# SEQN allows us to collect such disparate information for participants.
head(adult$SEQN)

# The following tells us that there is only one row in adult.csv for
# subject number 4.
dim(adult[adult$SEQN == 4,])

# Confirm the range stated on page 66 of the code book.
range(adult$SEQN)

# Confirm the race distribution stated on page 67 of the code book.
table(adult$DMARACER)
\end{verbatim}

\EXERCISE Continue using the adult survey data set in the following.

\hspace{10pt} a) Confirm the sex distribution stated on page 68 of the code book.

\hspace{10pt} b) Confirm the family size distribution using \texttt{HSFSIZER} (page 70).

\hspace{10pt} c) 
 Make a histogram of the age distribution.\vspace{-6pt}
\begin{verbatim}
hist(adult$HSAGEIR, breaks=seq(15,95,by=5), xlab="Age", 
     ylab="Percent per Year", probability=TRUE)
\end{verbatim}\vspace{-4pt}

\hspace{10pt} d) Confirm the census region distribution using \texttt{DMPCREGN} (page 71)
then make a barplot.\vspace{-6pt}
\begin{verbatim}
barplot(table(adult$DMPCREGN), names.arg=c("Northeast", "Midwest", "South", "West"))
\end{verbatim}\vspace{-4pt}

\hspace{10pt} e) Were the days of the week of the survey exams fairly uniform?
Use variable \texttt{MXPTIDW}.  If not, can you suggest a reason?

\hspace{10pt} f) What percent of respondents worked for the federal government?
Use variable HFD12 (page 124).  

\hspace{10pt} g) What percent of respondents worked for the either federal, state or local
 government?
Again, use variable HFD12 (page 124).  

\hspace{10pt} h) Did more respondents own cats or fish?  Use variables HFE8B and HFE8D.

\hspace{10pt} i) Is it possible that  more respondents owned either a bird or fish 
(or both) than  owned cats?  Use variables HFE8B and HFE8D.
\medskip

\EXERCISE  The variables on pages 131 -- 132 give information about smoking habits in 
 subjects' homes.  Let's construct a barplot of the number of cigarettes a day
smoked by person 1 in subjects' homes.  Do the following in RStudio.
\begin{verbatim}
cigs1  = adult$HFF3A
cigs1b = cigs1[cigs1 < 800]  # What are responses 888 and 999?
table(cigs1b)                # What does the value 777 mean?

cigs1b[cigs1b == 777] <- 0   # Look at the table below to see what this does.
table(cigs1b)

barplot(table(cigs1b))       # Now make the barplot.
\end{verbatim}
\vfill
\eject

\EXERCISE \textbf{Experimenting with the Exam Data}.  Open the \texttt{exam-acc.pdf}
code book.  In RStudio, load the exam data set.\vspace{-6pt}
\begin{verbatim}
exam = read.csv("/Users/username/Desktop/math207/exam.csv", header=TRUE)

# It's a data frame too.
class(exam) 

# How big is it?
dim(exam)

# Subject ID numbers are stored in both files.
head(adult$SEQN)
head(exam$SEQN)

# Here is survey information for subject 11.
adult[adult$SEQN == 11, ]

# Here is exam information for that subject.
exam[exam$SEQN == 11, ]

\end{verbatim}


\begin{verbatim}
# How tall (in cm) was subject 11? See page 199 of the code book
exam[exam$SEQN==11, ]$BMPHT

# How much did that subject weigh (in kg)?  See page 198.
exam[exam$SEQN==11, ]$BMPWT

# Let's examine the variation of weight with height.  
# Try the following then explain what goes wrong.
plot(exam$BMPHT, exam$BMPWT, xlab="height (cm)", ylab="weight (kg)")

# We need to extract the blank values.
w1 = exam[(exam$BMPWT != 888888) & (exam$BMPHT != 88888), ]$BMPWT
h1 = exam[(exam$BMPWT != 888888) & (exam$BMPHT != 88888), ]$BMPHT

plot(h1, w1, xlab="height (cm)", ylab="weight (kg)")
\end{verbatim}

Does weight increase with height?  Is the variation linear?  That is, does a straight
line fit the data  well?
\bigskip\bigskip\bigskip

\vfill
\eject
\end{document}


A (random) \textbf{variable} is a measurement that depends on the outcome 
   of a (random) event.
\begin{itemize}
\item \textbf{Quantitative} variables have numeric values.  
  \begin{itemize} 
     \item \textbf{Continuous} variables can assume a continuum of values:  Examples include 
      income, temperature, pressure, mass, and speed. 
     \item A \textbf{discrete} variable can assume only finitely (or countably) many values.
       Examples include: family size, and number of engine cylinders.
  \end{itemize}
\item \textbf{Qualitative} variables are non-numeric.  
   \begin{itemize}
      \item A qualitative variable can be \textbf{ordered} such as customer 
          satisfaction (good, better or best)
      \item or \textbf{unordered} such as eye color or automobile transmission type.
   \end{itemize}
\end{itemize}

We can compute averages, standard deviations and other numeric statistics for 
quantitative variables.  We can't do this for qualitative variables.

\begin{center}
{\footnotesize\begin{pspicture}(-1.5,0)(7.5,2.2)
\psset{linewidth=0.02, nodesep=2pt, xunit=3}
\rput(1,2){\rnode{Variables}{Variables}}
\rput(0,1){\rnode{Qualitative}{Qualitative}}
\rput(2,1){\rnode{Quantitative}{Quantitative}}
\rput(1.5,0){\rnode{Discrete}{Discrete}}
\rput(2.5,0){\rnode{Continuous}{Continuous}}
\rput(-0.5,0){\rnode{Ordered}{Ordered}}
\rput(0.5,0){\rnode{Unordered}{Unordered}}
\ncline{Variables}{Qualitative}
\ncline{Variables}{Quantitative}
\ncline{Quantitative}{Discrete}
\ncline{Quantitative}{Continuous}
\ncline{Qualitative}{Ordered}
\ncline{Qualitative}{Unordered}
\end{pspicture}}
\end{center}

\EXERCISE  Classify each of the variables encountered in the previous exercises
as discrete, continuous, ordered or unordered.

\vfill
\eject


\end{document}

